\scanpage{22}\mspage{7}
\uline{with universal curve} (\uline{in nonhomogeneous coordinates})
\[
\mathrm{Y}^{2}=4\,\mathrm{X}^{3}-\mathrm{g}_{2}\,\mathrm{X}-\mathrm{g}_{3}
\]
\margin{\handwritten{\uline{unit \unsure{section}} ?}}
\uline{and invariant differential} $\pi=\mathrm{d}\,\mathrm{X}/\mathrm{Y}$.

Moreover, when $2$ and $3$ are invertible, the datum of an invariant differential form suffices to rigidify a curve of genus $1$.

It follows immediately from (2.2) that we have
\[
\mathrm{c}_{4}^{3}-\mathrm{c}_{6}^{2}=\underset{\inserted{\handwritten{12^{3}}}}{1728}\,\Delta
\tag*{(2.6)}
\]
and this formula, reducing to an algebraic identity, is valid in every characteristic.

3.\qquad \uline{Case where} $2$ \uline{and} $\mathrm{c}_{4}$ \uline{are invertible}.

Suppose that $\mathrm{E}$ is a curve of genus one over a base $\mathrm{S}$ where $2$ is invertible. The invariant differential form $\pi$ being given, one can choose $\mathrm{x}$ and $\mathrm{y}$ such that in (1.3) one has $\mathrm{a}_{1}=\mathrm{a}_{3}=0$. The changes of coordinates preserving this condition are obtained by setting $\mathrm{u}=1$, $\mathrm{s}=\mathrm{t}=0$ in (1.2).

We thus have
\[
\mathrm{y}^{2}=\mathrm{x}^{3}+\mathrm{a}_{2}\,\mathrm{x}^{2}+\mathrm{a}_{4}\,\mathrm{x}+\mathrm{a}_{6}\,\raisebox{.14ex}{\scriptsize\unsure{,|'}}
\tag*{(3.1)}
\]
\[
\mathrm{b}_{2}=4\mathrm{a}_{2}\qquad\mathrm{b}_{4}=2\mathrm{a}_{4}\qquad
\mathrm{b}_{6}=4\mathrm{a}_{6}\qquad\mathrm{d}_{6}=8(\mathrm{a}_{2}\,\mathrm{a}_{4}-9\,\mathrm{a}_{6})
\]
\[
\mathrm{c}_{4}=2^{4}(\mathrm{a}_{2}^{2}-3\mathrm{a}_{4})\qquad
-\mathrm{c}_{6}=2^{6}\,\mathrm{a}_{2}^{3}-2^{5}\,3^{2}\,\mathrm{a}_{2}\,\mathrm{a}_{4}+2^{5}\,3^{3}\,\mathrm{a}_{6}\quad\text{and}
\tag*{(3.2)}
\]
\[
\Delta=2^{4}\,\delta\qquad,
\]
\clearpage
% archive sheet number: 21
